% Part: exercises
% Chapter: free logics
% Section: semantics

\documentclass[../../../../include/open-logic-section]{subfiles}

\begin{document}

\olfileid{exe}{frl}{sem}

\olsection{Semantics}

\begin{prob}[Semantics]
Explain why the following are not valid in any of the free logics.
(For neutral logic, feel free to pick just one trivalent logic.)
\begin{enumerate}
\item $\exists x(a=x)$
\item $\exists x(x=x)$
\end{enumerate}
\end{prob}

\begin{prob}[existence]
Explain why no formula of the form $\exists x\phi$ is valid in any
of the (inclusive) free logics. (For neutral free logic assume a truth-preservation
notion of validity.)
\end{prob}

\begin{prob}[Validity in negative and positive free logics]
For each of the formulas below, say whether it is valid in negative
free logic and whether it is valid in positive free logic. Briefly justify your answers.
\begin{enumerate}
\item $a=a$.\\
\emph{Example answer}. Positive: yes. Negative: no. In positive FL,
$t_{1}=t_{2}$ is true as long as $t_{1}$ and $t_{2}$ have the same
denotation, so $a=a$ is guaranteed to hold. In negative FL, $t_{1}=t_{2}$
is false if the denotation of $t_{1}$ or $t_{2}$ is in the outer
domain, so $a=a$ may fail. 
\item $\forall x(\lfrexists x)$
\item $Fa\rightarrow \lfrexists a$.
\item $\sim a=b\rightarrow\sim \lfrexists a$.
\item $\sim \lfrexists a\rightarrow\sim a=b$.
\item $\forall xGxb\rightarrow(\lfrexists a\rightarrow Gab)$.
\item $a=b\rightarrow(Fa\rightarrow Fb)$.
\item $\lfrexists a\land\sim \lfrexists b\rightarrow\sim Gab$.
\end{enumerate}
\end{prob}

\end{document}